\section{Introduction}\label{sec:introduction}
A learned continuation is useful in an economic decision only through the
policy it helps to implement. An accurate value prediction at familiar
states need not bound the payoff from all feasible deviations. A cheap
neural query need not repay the cost of learning, checking, and refreshing
the future that it evaluates. Neural Bellman Operators (NBO) address these
two questions together: they evaluate a specified future policy, improve
feasible current decisions, and attach an economic error and work account
to the resulting computation.

The paper retains a general controlled-economy formulation, including
recursive utility, endogenous preferences, temporal selves, and strategic
interaction. Its organizing object remains the policy-specific scalar
continuation. The continuation equation, the action actually implemented,
and the guarantee about its economic payoff are distinct. In a game,
evaluation against specified rivals is not an equilibrium certificate; with
recursive utility, a valid utility domain cannot be discarded when an
approximation is reused. These distinctions are essential to the common
architecture rather than exceptions to it.

The first part of the argument connects approximation to decisions.
Actionwise centered continuation errors control decision loss without
charging irrelevant level errors. Economic transport bounds then determine
when changes in transitions, rewards, or recursive evaluation permit reuse
and when they require refresh. A separate composition account propagates
full-action local advantages through the implemented policy's own future.
It displays state coverage, class approximation, value transfer, and
implementation separately. These premises cannot be supplied by a mean
loss calculated on a finite task catalogue.

This revision completes an all-state policy account in a multisector
capital-adjustment economy. A coercive Bellman-residual bound compares the
returned policy with every adapted finite-cost policy on an unbounded state
space. Its numerical instantiation has continuous correlated Gaussian
innovations, twenty-four dates, and ten or fifty independently adjustable
investment controls. All future policies are reoptimized in each regime.
The certificate evaluates the candidate's own quadratic value, encloses its
full vector-action improvement opportunity, and includes a declared
execution-error allowance. It does not use the optimal policy as an input
to the verifier, an action grid, or an extrapolation from visited states.

The neural implementation is also explicit. Every hidden weight of a
square-activation continuation is trained against its own policy's
population derivative-regression target. A trainable-factor proposition
connects the gradient and a verified rank bound, evaluated under the
returned policy, to its full-policy loss. This is a statement about a
jointly trained factor, not a fixed-feature least-squares assertion silently
applied to an arbitrary network. It does not assume that optimization
succeeds. The first successful outward economic check determines stopping;
failed checks and their candidates remain in the record.

The empirical comparison is deliberately strong. Conventional quadratic
continuations receive exactly the same own-policy evaluation information
and may take their closed-form minimizer. Cold and warm starts are compared
symmetrically for both representations. A separately timed Riccati solver
exploits the model's structure and is evaluated by the same verifier. The
neural and conventional representations span the same quadratic class,
while exact population evaluation is available in this economy. Thus this
experiment identifies the reliability and cost of a trained factor and
closes a previously conditional policy guarantee; it is not designed to
make simulation artificially expensive or to establish neural necessity.

The design fixes the complete methods, budgets, models, and a finite
initialization randomizer before primary execution. Sixteen initialization
seeds are enumerated, rather than treated as a sample from an unspecified
population. Success probabilities and expected capped operation counts
refer exactly to that declared randomizer. Wall clocks describe the
executed services and do not become population expectations through this
enumeration. Regimes change future technology or future stock valuation
while leaving the current primitives fixed. First-period investment
responses are enclosed using complete-policy loss bounds and refer to
fully reoptimized futures.

The earlier nonlinear capital experiments retain a different role. They
test continuation reuse, transport, scalar-interval welfare certification,
and complete first-success service costs under finite innovation laws.
Their adverse results remain visible. In the changing-future experiment,
adaptive NBO certifies all 120 of its future-specific outcomes, but costs
more than fresh NBO in every one of the eight dimension--volume cells;
cached sample-average optimization is the least-cost fully certified
procedure in those cells. Full-support enumeration also gives sharper
response bands. These findings establish accurate NBO decisions, not a
neural work advantage. Earlier protected payoff gains, cheaper Raw clocks,
unresolved comparisons, and nonpositive mechanism bounds all retain their
original targets and chronology.

There are therefore three kinds of conclusion in the paper. A verifier
result states what a certificate proves, independently of its candidate
generator. A candidate result establishes the accuracy of a particular
returned policy. A method result compares reliability or complete work
under a specified design. Neither of the first two implies the third. The
new full-policy experiment and the inherited nonlinear evidence are read
with this distinction throughout. The quadratic example does not discharge
the nonlinear utility-domain, information, and strategic-deviation
obligations of the other applications, and zero discretization error in a
declared discrete-time model is not a zero-error diffusion transfer.

The main article presents the economic targets, the continuation-to-policy
argument, the original NBO construction, and the all-state experiment. The
Prior Numerical Evidence companion contains the previously main-text
nonlinear comparisons in full; the technical supplement retains their
proofs and detailed records and adds the new arithmetic account. The
Economic Applications companion is preserved. This arrangement brings the
central proof and experiment together without deleting the historical
scientific content or changing the subject of the paper.
